%\documentclass[12pt,oneside]{amsart}
\documentclass{scrartcl}

%\documentclass{article}

\usepackage{amsthm,amsmath,amssymb}
\usepackage[numbers]{natbib}

%\usepackage[colorlinks]{hyperref}
\usepackage{hyperref}
\hypersetup{
    colorlinks,%
    citecolor=black,%
    filecolor=black,%
    linkcolor=black,%
    urlcolor=black
}
\usepackage{hypernat}
\usepackage[top=2.5cm, bottom=2.5cm, left=2.8cm,
right=2.8cm]{geometry}
\usepackage{float}

%\setlength{\parskip}{0pt}
%\setlength{\parsep}{0pt}
%\setlength{\headsep}{0pt}
%\setlength{\topskip}{0pt}
%\setlength{\topmargin}{0pt}
%\setlength{\topsep}{0pt}
%\setlength{\partopsep}{0pt}

\linespread{1}

%\usepackage{marvosym}

%\textwidth = 6.5 in \textheight = 9.2 in \oddsidemargin = 0.0 in
%\evensidemargin = 0.0 in \topmargin = -0.0 in \headheight = 0.0 in
%\headsep = 0.0in \parskip = 0.0in \parindent = 0.0in
%\def\baselinestretch{0.871}


\newtheorem{theorem}{Theorem}[section]
\newtheorem{proposition}[theorem]{Proposition}
\newtheorem{problem}[theorem]{Problem}
\newtheorem{fact}[theorem]{Fact}
\newtheorem{lemma}[theorem]{Lemma}
\newtheorem{defn}[theorem]{Definition}
\newtheorem{corollary}[theorem]{Corollary}
\newtheorem{remark}{Remark}[section]
\newtheorem{example}[theorem]{Example}

\newcommand{\E}{{\mathbb E}}
\newcommand{\se}{{\mathcal{E}}}
\newcommand{\R}{{\mathbb R}}
\renewcommand{\P}{{\mathbb P}}
\newcommand{\ac}{{\mathcal{A}}}
\newcommand{\A}{{\mathcal{A}}}
\newcommand{\B}{{\mathcal{B}}}
\newcommand{\C}{{\mathcal{C}}}
\newcommand{\M}{{\mathcal{M}}}
\newcommand{\Mf}{{\mathfrak{M}}}
\newcommand{\T}{{\mathcal{T}}}
\newcommand{\Z}{{\mathcal{Z}}}
\newcommand{\K}{{\mathcal{K}}}
\newcommand{\lc}{{\mathcal{L}}}
\newcommand{\Ps}{{\mathcal{P}}}
\newcommand{\G}{{\mathcal{G}}}
\newcommand{\pa}{{\mathring{p}}}
\newcommand{\F}{{\cal F}}
\newcommand{\samp}{{\mathcal{X}}}
\newcommand{\X}{{\mathcal{X}}}
\newcommand{\hi}{{\hat{\Phi}}}
\newcommand{\data}{{\text{data}}}

\newcounter{rcnt}[section]
\renewcommand{\thercnt}{(\roman{rcnt})}

\def\qt#1{\qquad\text{#1}}

\def\argmin{\mathop{\text{argmin}}}
\def\argmax{\mathop{\text{argmax}}}


\setlength{\parskip}{1.4 \medskipamount}
%\sloppy
%\linespread{1.3}

\begin{document}

\title{STAT 238 - Bayesian Statistics \\
 Homework Four} 
\subtitle{Spring 2026, UC Berkeley}
\author{Due by 11:59 pm on 03 April 2026 \\ Total Points = 118}

\date{}
\maketitle



% \tableofcontents

\begin{enumerate}

\item Consider the \texttt{spam7.csv} dataset which consists of data
  on $n = 4601$ emails corresponding to the following variables:
  \begin{itemize}
    \item \texttt{crl.tot}: Total length of uninterrupted sequences of capital letters.

    \item \texttt{dollar}: Occurrences of the dollar sign (\$), expressed as a percentage of the total number of characters.

    \item \texttt{bang}: Occurrences of the symbol \texttt{!}, expressed as a percentage of the total number of characters.

    \item \texttt{money}: Occurrences of the word ``money'', expressed as a percentage of the total number of words.

    \item \texttt{n000}: Occurrences of the string ``000'', expressed as a percentage of the total number of words.

    \item \texttt{make}: Occurrences of the word ``make'', expressed as a percentage of the total number of words.

    \item \texttt{yesno}: Indicates whether the message is spam, with levels
    \[
    n = \text{not spam}, \qquad y = \text{spam}.
    \]
  \end{itemize}
Our goal is to perform logistic regression on this dataset without
using any existing libraries. Our response variable $y$ will be the
\texttt{yesno} 
  variablle which indicates whether the email is spam or not. As
  covariates, use the following five variables
  \begin{align*}
&    x_1 = \log \left(\texttt{crl.tot} \right)  ~~~ x_2 =
                   \log(\texttt{dollar} + 0.001)  ~~~ x_3= \log(\texttt{bang} + 0.001)  \\
    & x_4 = \log \left(\texttt{money} + 0.001 \right) ~~~ x_5 = \log
      \left(\texttt{n000} + 0.001 \right) ~~~ x_6 = \log
      \left(\texttt{make} + 0.001 \right)
  \end{align*}
We are taking logarithms because these variables are highly skewed
otherwise. Also we are adding 0.001 because some of these variables
can take the value 0 and we don't want $\log (0)$.
\begin{enumerate}
\item Using the Newton's algorithm, calculate the maximum likelihood
  estimate of the coefficients in logistic regression. Check the
  correctness of your solution using an inbuilt library function for
  fitting logistic regression. (\textbf{4 points})

\item Compute the Laplace posterior normal approximation corresponding
  to the prior $\beta_0, \dots, \beta_m
  \overset{\text{i.i.d}}{\sim} \text{uniform}(-\infty,
  \infty)$. Compute the posterior standard errors for the
  coefficients. Check the correctness of your answers using an
  inbuilt library function for  fitting logistic
  regression. (\textbf{4 points})  
\end{enumerate}

 \item The \textbf{probit} regression model is given by the following. We
  observe data $(y_1, x_1), \dots, (y_n, x_n)$ with
  $y_1, \dots, y_n$ binary and $x_1, \dots, x_n \in \R^p$. $x_1,
  \dots, x_n$ are treated non-random and $y_1, \dots, y_n$ random. The
  probit regression model assumes 
  that $y_1 \dots, y_n$ are independent with
  \begin{equation*}
    y_i \sim \text{Bernoulli} \left(\Phi(x_i'
      \beta) \right)
  \end{equation*}
  where $\Phi(\cdot)$ is the standard normal cdf. The only difference
  from the logistic regression model is that $\exp(x_i'\beta)/(1 +
  \exp(x_i'\beta))$ is replaced by $\Phi(x_i'\beta)$. 
  \begin{enumerate}
  \item Write down the log-likelihood for this model and calculate its
    gradient and Hessian for a generic parameter vector $\beta$. (\textbf{3 +
      3 + 3 = 9 points}).
   \item Let $\hat{\beta}$ be the MLE. Consider the
  prior: $\beta_j \overset{\text{i.i.d}}{\sim} \text{uniform}(-\infty,
  \infty)$. Express the posterior normal
     approximation in terms of $\hat{\beta}$. (\textbf{3 points}). 
   \item Explicitly calculate the Laplace posterior normal approximation for
     the \textit{frogs} used in class in Lecture 17. Use the same response and
     covariate variables as in class. You can use standard library
     functions to calculate the MLE
     $\hat{\beta}$ or use Newton's method directly.  (\textbf{3 points})
\item Compare the posterior standard
     deviations (as calculated from the posterior normal
     approximation) to the standard errors reported by a library
     function performing standard frequentist
     analysis of probit regression. Are they similar?
     (\textbf{3  points}).
     \end{enumerate}

\item Consider the linear regression model:
  \begin{equation*}
    y_i = \beta_0 + \beta_1 x_i + \epsilon_i 
  \end{equation*}
for $i = 1, \dots, n$ with $\epsilon_i \overset{\text{i.i.d}}{\sim}
\text{Lap}(0, \sigma)$. Recall that $\text{Lap}(0, \sigma)$ is the
Laplace or double exponential density:
\begin{equation*}
x \mapsto  \frac{1}{2 \sigma} \exp \left(-\frac{|x|}{\sigma} \right)
\end{equation*}
Consider the usual prior:
\begin{equation*}
  \beta_0, \beta_1, \log \sigma \overset{\text{i.i.d}}{\sim}
  \text{Unif}(-\infty, \infty). 
\end{equation*}
The goal is to fit this model to $(x_1, y_1),
\dots, (x_n, y_n)$. 
\begin{enumerate}
\item Prove that   (\textbf{3 points})
  \begin{equation}\label{jden}
    f_{\beta_0, \beta_1 \mid \text{data}}(\beta_0, \beta_1) \propto
    \left(\frac{1}{S(\beta_0, \beta_1)} \right)^n
  \end{equation}
  where
  \begin{equation*}
    S(\beta_0, \beta_1) := \sum_{i=1}^n |y_i - \beta_0 - \beta_1
    x_i|. 
  \end{equation*}
 \item Consider the Pearson  father-son  height dataset that we used
   in class in Lecture 16. Fit this model to the data and report point estimates
   and standard errors of $\beta_0$ and $\beta_1$. You can approximate
   the joint density \eqref{jden} by the discrete distribution over a
   dense grid of $\beta_0$ and $\beta_1$. (\textbf{5 points})
  \item Compare your estimates of the intercept and slope with the
    corresponding estimates when fitting the usual regression model
    with normal errors. Comment on any discrepancies (or similarities)
    between the two sets of results. (\textbf{2 points}). 
\end{enumerate}


\item Consider the wages dataset $(x_i, y_i), i = 1, \dots, n$ from
  Lectures 18-21. $x_i$ denotes the number of years of experience, and
  $y_i$ denotes logarithm of weekly earnings for the $i$-th individual
  in the dataset. Suppose we want to fit the model:
  \begin{align*}
    y_i = \beta_0 + \beta_1 x_i + \beta_2 (x_i - c)_+ + \epsilon_i 
  \end{align*}
  with $\epsilon_i \overset{\text{i.i.d}}{\sim} N(0, \sigma^2)$. The
  unknown parameters in this model are $\beta_0, \beta_1, \beta_2,
  \sigma$ as well as $c$. The goal of this problem is to perform
  Bayesian inference for these  unknown parameters. Consider the
  prior: $\beta_0, \beta_1, \beta_2, c, \log \sigma$ are all
  independent with:
  \begin{align*}
    \beta_0, \beta_1, \beta_2, \log \sigma
    \overset{\text{i.i.d}}{\sim} \text{uniform}(-\infty, \infty)
  \end{align*}
  and
  \begin{align*}
    c \sim \text{uniform}(2, 60). 
  \end{align*}
  Note that $x_i$ ranges from $0$ to $63$, we are ruling out $c$
  taking values that are too close to the extreme points of the
  range. Other than this, the prior is uniformative about the value of
  $c$.

  \begin{enumerate}
  \item Show that the joint posterior density of $\beta, c, \sigma$ (here
    $\beta$ denotes the column vector with components $\beta_0,
    \beta_1, \beta_2$) is given by:
    \begin{align}\label{jopo}
      \text{posterior}(\beta, c, \sigma) \propto \sigma^{-n-1} \exp
      \left(-\frac{S(\beta, c)}{2 \sigma^2} \right) I\{\sigma > 0\}
      I\{2 < c < 60\}
    \end{align}
    where (\textbf{2 points})
    \begin{align*}
      S(\beta, c) = \sum_{i=1}^n \left(y_i - \beta_0 - \beta_1 x_i -
      \beta_2 (x_i - c)_+ \right)^2. 
    \end{align*}
  \item Let $X_c$ be the $n \times 3$ matrix whose $i$-th row is $1,
    x_i, (x_i - c)_+$ for $i = 1, \dots, n$. Prove that
    \begin{align}\label{pyth}
      S(\beta, c) = S(\hat{\beta}(c), c) + \left(\beta -
      \hat{\beta}(c) \right)^T X_c^T X_c  \left(\beta - \hat{\beta}(c)
      \right) 
    \end{align}
    where (\textbf{2 points})
    \begin{align*}
      \hat{\beta}(c) = \argmin_{\beta} \sum_{i=1}^n \left(y_i -
      \beta_0 - \beta_1 x_i - 
      \beta_2 (x_i - c)_+ \right)^2  = (X_c^T X_c)^{-1} X_c^T y. 
    \end{align*}
  \item Integrating $\beta$ from \eqref{jopo} (and using \eqref{pyth}
    to evaluate the integral), prove that the posterior of $c, \sigma$
    is: 
    \begin{align*}
      \text{posterior}(c, \sigma) \propto I\{\sigma > 0\}
      I\{2 < c < 60\} \sigma^{-n+p-1} |X_c^T X_c|^{-1/2} \exp
      \left(-\frac{S(\hat{\beta}(c), c)}{2 \sigma^2} \right)
    \end{align*}
    where $p = 3$ is the number of columns in $X_c$ and $|X_c^T X_c|$
    is the determinant of $X_c^T X_c$. (\textbf{3
      points})
  \item Integrate the above posterior over $\sigma$ to prove that the
    posterior of $c$ is: (\textbf{3 points})
    \begin{align*}
      \text{posterior}(c) \propto I\{2 < c < 60\} |X_c^T X_c|^{-1/2}
      \left(\frac{1}{S(\hat{\beta}(c), c)} \right)^{(n-p)/2}
    \end{align*}
  \item Numerically evaluate the posterior of $c$ for the wages
    data. Take a grid of values of $c$ in $(2, 60)$, compute the
    unnormalized posterior and then normalize it. Using the discrete
    approximation, compute and report a 95\% uncertainty interval for
    $c$. (\textbf{4 points}).
  \item Draw $N = 100$ posterior samples $(\beta^{(j)}_0, \beta^{(j)}_1,
    \beta^{(j)}_2, c^{(j)}, \sigma^{(j)})$ for $j = 1, \dots, N$. On
    the scatter plot of the data $(x_i, y_i), 1 \leq i \leq n$, plot
    each of the curves:
    \begin{align*}
      x \mapsto \beta^{(j)}_0 + \beta^{(j)}_1 x + \beta^{(j)}_2 (x -
      c^{(j)})_+. 
    \end{align*}
    Comment on the range of uncertainty revealed by the
    curves. (\textbf{4 points}). 
  \end{enumerate}

\item Consider \texttt{GolfTrends13March2026.csv} which contains
  monthly Google trends data for the query \textit{golf}. This is a
  monthly time series dataset that indicates the search popularity of
  \textit{golf} in the United States from January 2004 to February
  2026. The goal of this problem is to fit a smooth trend 
 function $\{\mu_t\}$ to the data $y_t$ using the model:
 \begin{equation*}
   y_t = \mu_t + \epsilon_t \qt{where $\epsilon_t
     \overset{\text{i.i.d}}{\sim} N(0, \sigma^2)$} 
 \end{equation*}
 and
 \begin{equation*}
   \mu_t = \beta_0 + \beta_1 (t-1) + \beta_2 (t - 2)_+ + \beta_3 (t -
   3)_+ + \dots + \beta_{n-1}(t - (n-1))_+. 
 \end{equation*}
 \begin{enumerate}
 \item A natural ridge regression estimator here would minimize:
\begin{equation*}
  \begin{split}
   & \sum_{t=1}^n \left(y_t  - \beta_0 - \beta_1 (t-1) - \beta_2 (t -
     2)_+ - \beta_3 (t - 3)_+ - \dots - \beta_{n-1}(t - (n-1))_+
     \right)^2 \\ &+ \lambda
      \sum_{j=2}^{n-1} \beta_j^2
  \end{split}
\end{equation*}  
  Implement the above (with a cross validation strategy for
   determining the regularization parameter) to estimate
   $\mu_t$. Report your value of $\lambda$ and plot the estimated
   trend function $\hat{\mu}_t$ on a plot of the data. Comment
   on whether the fitted trend function is capturing well the patterns
   present in the data without fully interpolating the
   data. (\textbf{6 points}) 
 \item Consider Bayesian inference with the prior:
   \begin{equation*}
     \beta_0, \beta_1\overset{\text{i.i.d}}{\sim} N(0, C), ~~ \beta_2,
     \dots, \beta_{n-1} 
     \overset{\text{i.i.d}}{\sim} N(0, \gamma^2 \sigma^2) 
   \end{equation*}
   along with the $\text{unif}(-C, C)$ prior for $\log(\sigma)$ and the
   $\text{unif}(\text{low}, \text{high})$ prior for $\log(\gamma)$ for
   appropriate $\text{low}$ and $\text{high}$. 

   \begin{enumerate}
   \item Generate $N = 1000$ posterior samples from $\gamma$ and
     $\sigma$. Based on these samples, calculate point estimates $\hat{\gamma}$ of
   $\gamma$ and $\hat{\sigma}$ of $\sigma$. How does your choice of the
   tuning parameter $\lambda$ in part (a) compare to
   $1/\hat{\gamma}^2$? (\textbf{4 points}).   
 \item Draw $N = 1000$ posterior samples from  $\{\mu_t\}$ and plot
  these sampled trend functions on a plot of the data. Comment
  on whether the uncertainty revealed in this plot makes
  sense. (\textbf{4 points})
\item Average over the $N$ posterior samples to obtain the Bayesian
  point estimate of $\hat{\mu}_t$. Compare this with the ridge
  regression estimate from part (a). Which trend estimate is smoother?
  Which trend estimate would you prefer? (\textbf{4 points}). 
\end{enumerate}
\end{enumerate}

\item Consider again the \textit{golf} google trends dataset as in the
  previous problem. In this problem, we will consider the model
  \begin{align*}
    y_t &= \beta_0 + \beta_1 (t-1) + \beta_2 (t - 2)_+ + \beta_3 (t -
   3)_+ + \dots + \beta_{n-1}(t - (n-1))_+ \\ &+ \beta_n \cos \left(2 \pi \frac{t}{12} \right) +
    \beta_{n+1} \sin \left(2 \pi \frac{t}{12} \right) + \epsilon_t. 
  \end{align*}
\begin{enumerate}  
 \item A natural ridge regression estimator here would minimize:
\begin{equation*}
  \begin{split}
    \sum_{t=1}^n \bigg(y_t & - \beta_0 - \beta_1 (t-1) - \beta_2 (t - 2)_+ - \beta_3 (t - 3)_+ - \dots - \beta_{n-1}(t - (n-1))_+ \\
    & - \beta_n \cos \left(2 \pi \frac{t}{12} \right) - \beta_{n+1}
      \sin \left(2 \pi \frac{t}{12} \right) \bigg)^2 + \lambda
      \sum_{j=2}^{n-1} \beta_j^2
  \end{split}                                                
\end{equation*}
   Implement this estimator with a cross validation strategy for
   determining the regularization parameter. Report your value of
   $\lambda$ and plot the estimated trend functions:
   \begin{equation*}
     \hat{\mu}_t := \hat{\beta}_0 + \hat{\beta}_1 (t-1) +
     \hat{\beta}_2 (t - 2)_+ + \hat{\beta}_3 (t - 
   3)_+ + \dots + \hat{\beta}_{n-1}(t - (n-1))_+
 \end{equation*}
 and
 \begin{equation*}
   \hat{\mu}_t + \hat{\beta}_n \cos \left(2 \pi \frac{t}{12} \right) + \hat{\beta}_{n+1}
      \sin \left(2 \pi \frac{t}{12} \right). 
 \end{equation*}
on a plot of the data. Comment
   on whether these trend functions capture well the patterns
   present in the data without fully interpolating the
   data. (\textbf{8 points}) 
 \item Consider Bayesian inference with the prior:
   \begin{equation*}
     \beta_0, \beta_1, \beta_n, \beta_{n+1}
     \overset{\text{i.i.d}}{\sim} N(0, C), ~~ \beta_2, \dots,
     \beta_{n-1} 
     \overset{\text{i.i.d}}{\sim} N(0, \sigma^2 \gamma^2)
   \end{equation*}
   along with the $\text{unif}(-C, C)$ prior for $\log(\sigma)$ and
   $\text{unif}(\text{low}, \text{high})$ prior for $\log(\gamma)$ for
   appropriately chosen $\text{low}$ and $\text{high}$ ($C \rightarrow
   \infty$). 
   \begin{enumerate}
   \item Generate $N = 1000$ posterior samples from $\gamma$ and
     $\sigma$. Using these samples, calculate point estimates $\hat{\gamma}$ of
   $\gamma$ and $\hat{\sigma}$ of $\sigma$. How does your choice of the
   tuning parameter $\lambda$ in part (a) compare to
   $1/\hat{\gamma}^2$? (\textbf{5 points}).   
 \item Draw $N = 1000$ posterior samples from  $\{\mu_t\}$ (note that
   $\mu_t$ does not include the sinusoidal part) and plot
  these sampled trend functions on a plot of the data. Comment
  on whether the uncertainty revealed in this plot makes
  sense. (\textbf{5 points})
\item Average over the $N$ posterior samples to obtain the Bayesian
  point estimate of $\hat{\mu}_t$. Compare this with the ridge
  regression estimate from part (a). Which trend estimate is smoother?
  Which trend estimate would you prefer? (\textbf{5 points}).
\item Obtain Bayesian posterior mean estimates of $\hat{\beta}_n$
  and $\hat{\beta}_{n+1}$. Plot the function:
  \begin{align*}
    \hat{\mu}_t + \hat{\beta}_n \cos \left(2 \pi \frac{t}{12} \right) + \hat{\beta}_{n+1}
      \sin \left(2 \pi \frac{t}{12} \right)
  \end{align*}
  on a plot of the data (here $\hat{\mu}_t$ is the posterior mean
  estimate from part (iii)). Comment on how well this function
  captures the underlying trend present in the data. (\textbf{4
    points}). 
\end{enumerate}
\end{enumerate}  
\item Consider the wages dataset $(x_i, y_i), i = 1, \dots, n$ from
  Lectures 18-21. $x_i$ denotes the number of years of experience, and
  $y_i$ denotes logarithm of weekly earnings for the $i$-th individual
  in the dataset. Because this dataset is large, take a random
  subsample of size $n = 500$ by setting the random seed of $42$ (as
  in Lecture 20). The goal of this problem is to fit the model:
  \begin{align*}
    y_i = f(x_i) + \epsilon_i 
  \end{align*}
  with
  \begin{align*}
    f(x_i)  = \beta_0 + \tau B(x_i)
  \end{align*}
  where $\beta_0 \sim \text{unif}(-C, C)$ and $B(\cdot)$ is Brownian
  motion. Also $\epsilon_i \overset{\text{i.i.d}}{\sim} N(0,
  \sigma^2)$. Write $\tau = \gamma \sigma$, and take the
  $\text{unif}(-C, C)$ prior on $\log \sigma$ as well as the prior
  $\text{unif}(\text{low}, \text{high})$ on $\log(\gamma)$ for two
  appropriate values $\text{low}$ and $\text{high}$.
  \begin{enumerate}
  \item Describe the prior distribution of $f$ as a Gaussian Process
    with the correct mean and covariance kernel. (\textbf{1 point})
  \item Describe the posterior distribution of $f$ given the data and
    given the values of $\sigma$ and $\tau$. (\textbf{3 points})
  \item Assume $\tau = \sigma \gamma$. Describe the posterior
    distribution of $\sigma$ given the data and given the value of
    $\gamma$. (\textbf{3 points})
  \item Describe the marginal posterior distribution of
    $\gamma$. (\textbf{3 points}).
  \item Implement this posterior inference on the data of size $500$
    and compute the posterior mean estimate of the function $f$ on a
    grid of values in the range $[0, 63]$. Plot this posterior mean on
    a plot of the data. Compare this estimate with that derived from
    the method of Lecture 20. (\textbf{4 points}).
  \item Repeat part (e) for the full dataset. (\textbf{5
      points}). 
  \end{enumerate}

\item \textbf{[Bonus Question]}: Implement any of the exercises in this
  homework using PyMC. If PyMC does not work as expected—for example,
  if the MCMC algorithm fails to converge, produces clearly incorrect
  results, or takes an unusually long time to run—please report the
  issue. You will receive 5 bonus points for each problem for which
  you submit a well-documented report of such an issue. 
\end{enumerate}


%\begin{thebibliography}{99}

  
%\bibitem[Fisher(1934)]{Fisher1934}
%R.~A. Fisher,
%\newblock ``Probability, likelihood and quantity of information in the logic of
%uncertain inference,''
%\newblock \emph{Proceedings of the Royal Society of London. Series~A,
%Containing Papers of a Mathematical and Physical Character}, vol.~146,
%no.~856, pp.~1--8, 1934.

%\bibitem[Efron(2010)]{Efron}
%Efron, B. (2010)
%\textit{Large-scale inference: empirical Bayes methods for estimation,
%testing, and prediction}. Cambridge University Press, 2010.

%\bibitem[Gelman et~al.(2013)]{Gelman2013}
%Gelman, A., Carlin, J. B., Stern, H. S., Dunson, D. B., Vehtari, A., and Rubin, D. B. (2013)
%\textit{Bayesian data analysis}. Third edition. Chapman \& Hall/CRC, 2013.

%     \bibitem[Jaynes(2003)]{Jaynes} Jaynes, E. T, (2003)
%  \textit{Probability theory: the logic of science}. Cambridge
%  University Press, 2003.

%  \bibitem[Sinai(1992)]{Sinai} Sinai, Y. G, (1992)
%  \textit{Probability theory: an introductory course}. Springer
%  Verlag, 1992.  

%\bibitem[{deGroot(1986)}]{DeGrootBlackwell}DeGroot, M. H. (1986)
%       A conversation with David Blackwell.
%\newblock \emph{Statistical Science}, \textbf{1}, 40--53.

%         \bibitem[Mosteller(1987)]{Mosteller} Mosteller, F, (1987)
%  \textit{Fifty challenging problems in probability with
%    solutions}. Courier Corporation, 1987.

%    \bibitem[MacKay(2003)]{MacKay} MacKay, D. J. C., (2003)
% \textit{Information theory, inference, and learning
%  algorithms}. Cambridge University Press, 2003.

%\bibitem[{Lindley(2013)}]{Lindley}Lindley, D. V. (2013)
%   \textit{Understanding Uncertainty}. John Wiley and sons, 2013.


%\end{thebibliography}


\end{document}

%%% Local Variables:
%%% mode: latex
%%% TeX-master: t
%%% End:
